
\clearpage
\appendices


% Appendix-style algorithm numbering (A1, A2, ...)
\setcounter{algorithm}{0}
\renewcommand{\thealgorithm}{A\arabic{algorithm}}

\newenvironment{appalgorithm}[2]{%
  \refstepcounter{algorithm}\label{#2}%
  \Needspace{10\baselineskip}% avoid splitting header from first lines
  \par\medskip
  \noindent\hrule height 0.4pt\relax
  \vspace{0.35em}
  \noindent\textbf{Algorithm \thealgorithm.} #1\par
  \vspace{0.35em}
  \small
}{%
  \par\normalsize
  \vspace{0.35em}
  \noindent\hrule height 0.4pt\relax
  \par\medskip
}

% ============================================================
% Appendix A
% ============================================================
\section{Benchmark Construction and Chunking Algorithms}
\label{app:benchmark_algorithms}

\noindent This appendix summarizes three procedures used in benchmark construction and context-bounded migration:
(i) relevance scoring, (ii) constraint-aware selection, and (iii) PHP-aware smart chunking.

% ------------------------------------------------------------
\subsection{Benchmark construction}
\label{app:benchmark_construction}

\paragraph{Relevance scoring.}
We score candidate files by contract-triggered rule activity and migration relevance signals.

\begin{appalgorithm}{File scoring for migration relevance under the pinned contract.}{alg:file_scoring}
\begin{algorithmic}[1]
  \Require file with $LOC$, $V$ (\texttt{php\_version\_changes}), and triggered rules $R$
  \Ensure relevance score $s$
  \State $s \gets 10 \cdot |R| + 5 \cdot V$
  \If{$100 \le LOC \le 2000$}
    \State $s \gets s + 20$
  \ElsIf{$LOC > 2000$}
    \State $s \gets s - 0.01 \cdot (LOC - 2000)$
  \EndIf
  \State $s \gets s + 8 \cdot \left|\{\text{PHP rule families covered by } R\}\right|$
  \State $s \gets s + \textsc{ImportantRuleBonus}(R)$
  \State \Return $s$
\end{algorithmic}
\end{appalgorithm}

\paragraph{Constraint-aware selection.}
We select 100 files while enforcing full rule coverage as the hard constraint.

\begin{appalgorithm}{Constraint-aware selection of 100 files with full rule coverage.}{alg:selection_engine}
\begin{algorithmic}[1]
  \Require summary table, rule metadata, size bins, target counts $T$
  \Ensure selected file set $S$
  \State group files by size category
  \State compute score for each file using Algorithm~\ref{alg:file_scoring}
  \ForAll{size categories $c$}
    \State select top $T[c]$ files by score into $S$
  \EndFor
  \State $S \gets \textsc{EnsureRuleCoverage}(S)$
    \Comment{May add files outside bin quotas; coverage is the hard constraint}
  \If{$|S| > 100$}
    \State trim lowest-scoring files if coverage is preserved
  \ElsIf{$|S| < 100$}
    \State add highest-scoring remaining files until $|S| = 100$
  \EndIf
  \State \Return $S$
\end{algorithmic}
\end{appalgorithm}

% ------------------------------------------------------------
\subsection{PHP-aware chunking for bounded-context migration}
\label{app:chunking}

\paragraph{Smart chunking.}
We split long PHP files into bounded-context segments while attempting to avoid mid-function splits under a fixed policy.

\begin{appalgorithm}{PHP-aware smart chunking.}{alg:smart_chunking}
\begin{algorithmic}[1]
\Require $code$ (PHP source string), $chunk\_size$ (target lines), $MAX\_EXTENSION$ (max extension in lines)
\Ensure $chunks$ (list of records $\{start\_line,end\_line,actual\_size,total\_lines,code\}$)
\State $lines \gets \Call{Split}{code,\ "\backslash n"}$
\State $total \gets |lines|$
\If{$total \le chunk\_size$}
  \State \Return $[\{1,total,total,total,code\}]$
\EndIf
\State $boundaries \gets \Call{FindFunctionBoundaries}{code}$
\Comment{Tokenizer-guided detection via \texttt{token\_get\_all} when available; regex fallback; end lines approximated via brace matching}
\State $chunks \gets [\ ]$
\State $current \gets 0$
\While{$current < total$}
  \State $initial\_end \gets \min(current + chunk\_size - 1,\ total - 1)$
  \State $final\_end \gets initial\_end$
  \State $relevant \gets \{f \in boundaries \mid (current \le f.start \le initial\_end)\ \lor\ (f.start < current \land f.end \ge current)\}$
  \ForAll{$f \in relevant$}
    \If{$f.end \le initial\_end + MAX\_EXTENSION$}
      \State $final\_end \gets \max(final\_end, f.end_prof)$
      \State $final\_end \gets \max(final\_end, f.end)$
    \EndIf
  \EndFor
  \State $chunk\_code \gets \Call{Join}{lines[current..final\_end],\ "\backslash n"}$
  \State append $\{current+1,final\_end+1,(final\_end-current+1),total,chunk\_code\}$ to $chunks$
  \State $current \gets final\_end + 1$
\EndWhile
\State \Return $chunks$
\end{algorithmic}
\end{appalgorithm}

% ============================================================
% Appendix B
% ============================================================

\section{Prompt Templates}
\label{app:prompt_templates}

We use prompting shared across models, with whole-file and chunk-segment variants.
In both templates, \texttt{\{code\}} is replaced with the input PHP text.

\subsection{Whole-file prompt}
\noindent\begin{Verbatim}[breaklines=true,breakanywhere=true,fontsize=\footnotesize]
You are a senior PHP developer with expertise in legacy code modernization.
Your task is to migrate this legacy PHP code up to PHP 8.3 standards using modern syntax and features while maintaining functional equivalence.

Please migrate the following PHP code up to PHP 8.3 standards:

{code}

Your response should follow this EXACT format:

// MIGRATION_START
[your migrated PHP code here]
// MIGRATION_END

CRITICAL FORMATTING REQUIREMENT:
- Place the MIGRATION_START marker BEFORE the opening <?php tag
- Place the MIGRATION_END marker AFTER the closing PHP code
- Do NOT place these markers inside the PHP code itself

Provide only the migrated PHP code with the markers placed correctly outside the PHP code block, no additional commentary.
\end{Verbatim}

\subsection{Chunk-segment prompt}
The placeholders \texttt{\{filename\}}, \texttt{\{start\_line\}}, \texttt{\{end\_line\}}, \texttt{\{total\_lines\}}, \texttt{\{chunk\_number\}}, and \texttt{\{total\_chunks\}} are filled from the chunk metadata.

\noindent\begin{Verbatim}[breaklines=true,breakanywhere=true,fontsize=\footnotesize]
You are a senior PHP developer with expertise in legacy code modernization.
Your task is to migrate this PARTIAL SEGMENT of a larger PHP file up to PHP 8.3 standards using modern syntax and features.

CONTEXT:
- Original file: {filename}
- Processing lines: {start_line} to {end_line} (of {total_lines} total lines)
- This is chunk {chunk_number} of {total_chunks}

CRITICAL INSTRUCTIONS FOR PARTIAL CODE SEGMENTS:
1. This is ONLY a SEGMENT of a larger file - DO NOT try to complete it
2. keep the original opening <?php tag if present. if the segment does not start with <?php, DO NOT add one
3. keep the original closing ?> tag if present. If the segment does not end with ?>, DO NOT add one
4. keep the original closing braces }} if present. DO NOT add any closing braces that are not in the original segment
5. keep the original opening braces {{ if present. DO NOT add any opening braces that are not in the original segment
6. DO NOT try to complete class definitions, function definitions, or any code structures
7. Preserve the EXACT START and END boundaries of the provided code segment

WARNING: Adding extra braces that are not in the original segment or completing code structures will break the reconstruction process!

Please modernize ONLY the following PHP code segment up to PHP 8.3 standards:

{code}

Your response should follow this EXACT format:

// MIGRATION_START
[your modernized code segment here - exactly as provided, no additions]
// MIGRATION_END

CRITICAL FORMATTING REQUIREMENT:
- Place the MIGRATION_START marker BEFORE the code segment
- Place the MIGRATION_END marker AFTER the code segment
- Do NOT place these markers inside the PHP code itself

Migrate only the provided code segment. Do not add any missing parts or try to complete incomplete structures.
\end{Verbatim}

% ============================================================
% Appendix C
% ============================================================
\section{Loadability Tripwire Slice Specification}
\label{app:runtime_validation}

This appendix documents the \emph{isolated, stubbed loadability tripwire} used on an execution slice to surface harness-visible load-time failures after migration.
It specifies the execution environment, harness bootstrap, slice accounting, outcome taxonomy, and diagnostic coverage.

\subsection{Scope and interpretation}
This signal is \emph{harness-dependent} and is not an integration test or an end-to-end correctness guarantee.
We interpret it primarily as a \emph{regression tripwire}: \emph{New failure} (Orig.\ Pass $\rightarrow$ Migr.\ Fail).
Unless otherwise stated, all percentages are computed over the full slice size $N_{\mathrm{rt}}=55$ (no files are excluded).
We therefore treat this tier as a \emph{within-harness regression detector} whose purpose is to flag newly introduced load-time failures under identical conditions.

\subsection{Regression definition and harness-scoped interpretation}
A \textbf{regression} is defined as a file that loaded successfully in the original snapshot (under PHP 8.3.22 + stubs) but fails to load after migration.
These are harness-visible include-time failures: parse errors, fatal errors, uncaught exceptions, and missing-symbol errors that prevent \texttt{require\_once} from completing.

Observed regression counts across the slice ($N_{\mathrm{rt}}=55$):

\begin{center}
\footnotesize
\begin{tabular}{@{}lrr@{}}
\toprule
\textbf{Model} & \textbf{Regressions} & \textbf{Regression Rate} \\
\midrule
Rector Baseline & 7 & 12.7\% \\
Claude Sonnet 4 & 8 & 14.5\% \\
GPT-5 Codex & 9 & 16.4\% \\
Llama 3.3 70B & 12 & 21.8\% \\
Gemini 2.5 Pro & 12 & 21.8\% \\
Gemini 2.5 Flash & 18 & 32.7\% \\
\bottomrule
\end{tabular}
\end{center}

\noindent\textbf{Interpretation scope.}
These failures represent \emph{harness-visible include-time defects}---i.e., the file cannot be included without a fatal condition under the tripwire harness---but they are detected under isolated-stub conditions, not full application bootstrap.
Accordingly, the harness does not validate behavioral correctness, API contracts, or end-to-end functionality; it detects only include-time failures.
A file that passes this tripwire may still contain behavioral regressions; conversely, some files that fail may succeed under full WordPress initialization.
We treat \emph{New failure} counts as a \emph{lower bound} on migration breakage under the harness: they surface include-time breaks (e.g., syntax errors, missing symbols, invalid includes, or signature/parameter drops that trigger immediate fatals) that static contract discharge may not detect, but do not capture semantic or integration-level issues.

%------------------------------------------------------------------------------
\subsection{Execution environment}
\label{appendix:exec-env}

All tripwire checks are performed under a fixed environment for repeatability.
\textbf{Execution environment:} We execute all code (original and migrated) under PHP 8.3.22.
When we refer to ``original PHP 7.4 code,'' this denotes the source snapshot intended for PHP 7.4 compatibility, not the interpreter version used for execution.

\begin{center}
\footnotesize
\captionof{table}{Execution Environment Specification (Loadability Tripwire)}
\label{tab:runtime-env}
\vspace{-0.3em}
\begin{tabular}{@{}ll@{}}
\toprule
\textbf{Component} & \textbf{Specification} \\
\midrule
PHP Version & 8.3.22 (cli) \\
Operating System & Windows 11 (x64) \\
Memory Limit & 512 MB \\
Execution Timeout & 30 seconds per file \\
Error Reporting & \texttt{E\_ALL} \\
\bottomrule
\end{tabular}
\end{center}

\noindent\emph{Note.}
We use a single fixed OS and interpreter configuration for repeatability; this tier is interpreted only via regression direction (Orig.\ Pass $\rightarrow$ Migr.\ Fail) under the same harness, not as a portable runtime correctness metric.
Because include-time behavior can be OS- and configuration-sensitive (e.g., filesystem case-sensitivity and path handling, extension availability), we treat cross-environment generalization as out of scope for this signal.

%------------------------------------------------------------------------------
\subsection{Execution slice definition ($N_{\mathrm{rt}}=55$)}
\label{appendix:runtime-slice-def}

We evaluate a fixed execution slice of size $N_{\mathrm{rt}}=55$ benchmark files.
Each file is executed in isolation (fresh PHP process) under a minimal WordPress-core stub environment and then required once.
We record \textbf{loadability outcomes} (load success vs.\ fatal error/exception).

The tripwire checks only whether each file loads without fatal errors; it does not validate behavioral correctness, API contracts, or end-to-end functionality.
Some files may fail to load even in the original snapshot due to isolation with stubs (e.g., external dependencies, cross-file coupling, or missing runtime initialization).
\textbf{We retain all 55 files in-slice} and represent such cases via the outcome taxonomy in Section~\ref{appendix:outcome-taxonomy} rather than excluding them.

%------------------------------------------------------------------------------
\subsection{Harness bootstrap}
\label{appendix:harness-bootstrap}

\subsubsection{File loading procedure}
Files are loaded using the following procedure implemented in \texttt{run\_one.php}:

\begin{enumerate}
    \item Load WordPress core stubs (\texttt{stubs/wp\_core.php})
    \item Load translation entry stubs (\texttt{stubs/entry.php})
    \item Include the file under test via \texttt{require\_once}
    \item Record load success or fatal error
\end{enumerate}

\subsubsection{Stubs and mocks}
The stub file \texttt{wp\_core.php} provides 82 function stubs, 7 class stubs, and 20 constant definitions.
Table~\ref{tab:stubs} summarizes the stubbed APIs and globals.
The harness does \emph{not} model full WordPress bootstrap; as a result, outcomes can reflect harness-visible interface drift, missing initialization, or unmodeled dependencies rather than end-to-end behavior.

% NOTE: This is the only wide table (table*) to avoid float pile-ups in 2-column layout.
\begin{table*}[t]
\centering
\footnotesize
\setlength{\tabcolsep}{5pt}
\renewcommand{\arraystretch}{1.15}
\caption{Stubbed WordPress APIs and Globals (Loadability Tripwire Harness)}
\label{tab:stubs}
\begin{tabularx}{\textwidth}{@{}>{\raggedright\arraybackslash}p{2.8cm}X@{}}
\toprule
\textbf{Category} & \textbf{Stubbed Elements} \\
\midrule
Global Variables &
\texttt{\$wpdb}, \texttt{\$wp\_filter}, \texttt{\$wp\_actions}, \texttt{\$wp\_current\_filter}, \texttt{\$wp\_query}, \texttt{\$post}, \texttt{\$wp\_rewrite} \\
I18n Functions &
\texttt{\_\_()}, \texttt{\_e()}, \texttt{\_x()}, \texttt{\_n()}, \texttt{\_nx()}, \texttt{esc\_html\_\_()}, \texttt{esc\_attr\_\_()} \\
Hook System &
\texttt{add\_action()}, \texttt{remove\_action()}, \texttt{do\_action()}, \texttt{add\_filter()}, \texttt{remove\_filter()}, \texttt{apply\_filters()}, \texttt{has\_filter()}, \texttt{has\_action()}, \texttt{did\_action()}, \texttt{current\_filter()} \\
Sanitization &
\texttt{esc\_url()}, \texttt{esc\_js()}, \texttt{esc\_sql()}, \texttt{wp\_kses\_post()} \\
Utilities &
\texttt{wp\_parse\_args()}, \texttt{absint()}, \texttt{wp\_unslash()}, \texttt{wp\_slash()}, \texttt{trailingslashit()} \\
User Functions &
\texttt{wp\_get\_current\_user()}, \texttt{get\_current\_user\_id()}, \texttt{current\_user\_can()} \\
Option/Transient &
\texttt{get\_option()}, \texttt{update\_option()}, \texttt{delete\_option()}, \texttt{get\_transient()}, \texttt{set\_transient()} \\
Conditional Tags &
\texttt{is\_admin()}, \texttt{is\_multisite()}, \texttt{is\_user\_logged\_in()}, \texttt{is\_singular()} \\
Classes &
\texttt{WP\_Error}, \texttt{WP\_Widget}, \texttt{WP\_List\_Table}, \texttt{WP\_Dependencies}, \texttt{IXR\_Client}, \texttt{ftp\_base}, \texttt{Translation\_Entry} \\
\midrule
\textit{Not modeled} &
Database connections, HTTP requests, filesystem writes, session state, full WordPress core bootstrap \\
\bottomrule
\end{tabularx}
\end{table*}

\subsubsection{Error handling policy}
The runner implements a custom error handling policy:

\begin{itemize}
    \item \textbf{Fatal errors / Exceptions}: immediate failure (captured via try-catch)
    \item \textbf{Errors (E\_ERROR, E\_PARSE)}: immediate failure
    \item \textbf{Warnings (E\_WARNING)}: recorded; treated as failure only when attributed to the target file (by \texttt{\$errfile} path prefix match)
    \item \textbf{Notices / Deprecations (E\_NOTICE, E\_DEPRECATED)}: recorded but tolerated; logged for analysis
\end{itemize}

%------------------------------------------------------------------------------
\subsection{Slice accounting}
\label{appendix:slice-accounting}

The harness is configured over a fixed slice of $N_{\mathrm{rt}}=55$ target files.
Some files have known isolation limitations (external deps / cross-file coupling) but remain included in the slice and denominator.

\begin{center}
\footnotesize
\captionsetup{skip=5pt, justification=raggedright, singlelinecheck=false}
\setlength{\tabcolsep}{4pt}
\renewcommand{\arraystretch}{1.15}
\captionof{table}{Files with Known Isolation Limitations (Included in Slice)}
\label{tab:isolation-limited}
\begin{tabularx}{\columnwidth}{@{}>{\raggedright\arraybackslash}p{0.42\columnwidth}
                                >{\raggedright\arraybackslash}X@{}}
\toprule
\textbf{File} & \textbf{Reason} \\
\midrule
\texttt{getid3.php} & Requires getID3 core library classes. \\
\texttt{module.audio.ac3.php} & Requires \texttt{getid3\_handler} base class. \\
\texttt{class-ftp.php} & Cross-file dependency on FTP classes. \\
\texttt{class-pop3.php} & Depends on network/socket initialization. \\
\texttt{class-json.php} & Requires \texttt{Services\_JSON} library. \\
\bottomrule
\end{tabularx}
\end{center}

%------------------------------------------------------------------------------
\subsection{Outcome taxonomy}
\label{appendix:outcome-taxonomy}

A file is counted as \textbf{Pass} for a given codebase version iff it loads successfully without fatal errors.
Include-time fatal errors (parse errors, fatal runtime errors, exceptions) are counted as failures.

\begin{center}
\footnotesize
\captionof{table}{Outcome Categories for the Isolated Loadability Tripwire}
\label{tab:runtime-outcomes}
\vspace{0.3em}
\begin{tabular}{@{}lp{0.72\columnwidth}@{}}
\toprule
\textbf{Outcome} & \textbf{Interpretation} \\
\midrule
Both pass & Both original and migrated load successfully. \\
Fixed failure & Original fails to load, but migrated loads successfully. \\
New failure & Original loads, but migrated fails to load. \\
Both fail & Both original and migrated fail to load. \\
\bottomrule
\end{tabular}
\end{center}

For concision in the main text, we aggregate:
\textbf{Pass} = (Both pass + Fixed failure), \textbf{Regress} = (New failure), and \textbf{Other} = (Both fail).

%------------------------------------------------------------------------------
\subsection{Harness sanity checks}
\label{appendix:bias-controls}

For transparency, we report absolute pass rates (load success) under identical harness conditions, computed over the full $N_{\mathrm{rt}}=55$ slice.
Because the harness does not boot WordPress and files are executed in isolation under stubs, these pass rates are harness-dependent and are \emph{not} interpreted as integration or correctness measures; we use this tier primarily to surface \emph{Regress} outcomes.

\begin{center}
\footnotesize
\captionof{table}{Harness Pass Rates (Context Only; Not a Correctness Metric; $N_{\mathrm{rt}}=55$)}
\label{tab:bias-control}
\vspace{0.3em}
\begin{tabular}{@{}lrr@{}}
\toprule
\textbf{Codebase Version} & \textbf{Passed} & \textbf{Pass Rate} \\
\midrule
Original snapshot  & 50 & 90.9\% \\
Rector-migrated baseline & 46 & 83.6\% \\
Claude Sonnet 4 & 44 & 80.0\% \\
GPT-5 Codex & 43 & 78.2\% \\
Gemini 2.5 Pro & 40 & 72.7\% \\
Llama 3.3 70B & 38 & 69.1\% \\
Gemini 2.5 Flash & 33 & 60.0\% \\
\bottomrule
\end{tabular}
\end{center}

\noindent We therefore emphasize \textbf{Regress} counts in the main text and do not use absolute pass rates for model ranking.

\noindent\textbf{Interpretation.}
The non-ceiling pass rate of the original snapshot (90.9\%) indicates unavoidable isolation limitations (Table~\ref{tab:isolation-limited}): 5 files fail to load due to external dependencies or deprecated PHP syntax (curly-brace array access) removed in PHP~8.
Such cases are represented explicitly via the \emph{Both fail / Other} outcome rather than excluded.

%------------------------------------------------------------------------------
\subsection{Failure exemplars}
\label{appendix:failure-exemplars}

Representative harness-visible load-time failures from actual test results include:

\begin{enumerate}[leftmargin=*, itemsep=2pt, topsep=2pt, parsep=0pt]
  \item \textbf{Deprecated curly-brace array access syntax (original files).}\\
  \emph{Files:} \texttt{class-ftp.php}, \texttt{class-pop3.php}, \texttt{getid3.php}, \texttt{module.audio.ac3.php}, \texttt{class-json.php}.
  Original files use \texttt{\$var\{0\}} syntax (removed in PHP~8) or require external dependencies, causing 5 of 55 original files to fail (90.9\% baseline).

  \item \textbf{Callback transformation drops required parameters (Rector).}\\
  \emph{Files:} \texttt{class.akismet-widget.php}, \texttt{user.php}, \texttt{file.php}.
  Rector converted closures to arrow functions and dropped required callback parameters, triggering fatal argument-count errors at callback invocation sites under the harness.

  \item \textbf{Parse error from malformed syntax (Gemini models).}\\
  \emph{Files:} \texttt{Parser.php}, \texttt{File.php}.
  Gemini 2.5 Flash introduced malformed syntax including unexpected tokens (``unexpected token ?'' or ``unexpected token (''), causing include-time parse failures.

  \item \textbf{Modified include path points to missing file (multiple models).}\\
  \emph{File:} \texttt{translations.php}.
  Both Rector and Gemini 2.5 Pro modified \texttt{require\_once} paths, causing ``Failed opening required'' errors when the referenced file doesn't exist.

  \item \textbf{High regression rate correlates with low resolution.}\\
  Gemini 2.5 Flash showed 18 regressions (32.7\% regression rate) with 53.5\% rule resolution, while Rector showed 7 regressions (12.7\%) with 98.3\% resolution, demonstrating the tripwire's sensitivity to incomplete migrations.
\end{enumerate}

\vspace{0.25em}
\noindent\textbf{Summary.}
The loadability tripwire detects include-time failures and harness-visible load-time regressions (parse errors, fatal errors, missing dependencies, broken inheritance) that static contract discharge may not surface.
Because all metrics are computed over a fixed $N_{\mathrm{rt}}=55$ slice, unavoidable isolation limitations are represented via the \emph{Both fail / Other} category rather than excluded.
This tier provides a complementary regression signal but does not validate behavioral correctness or end-to-end functionality.
